\chapter{The Italian Silver Economy and the Cost of an Unsegmented Market}
\label{ch:ch1}

\section{A demographic transformation with an economic price tag}
Italy is ageing faster than almost any other country in the European Union and at a pace that decisively reframes its over-65 population from a fiscal and clinical liability into the largest single addressable consumer segment of the next two decades. According to Istat (2024), roughly 14.2~million Italian residents were aged 65 or older at the start of 2024, equivalent to 23.8\% of the total population and the second-highest share among the EU-27 after Germany. Eurostat projects this share to reach 30\% by 2050, with the old-age dependency ratio rising from 37.5 to 60. The European Commission's silver-economy reports (European Commission, 2018, 2021) estimate that the EU's 50-plus segment already generates a quarter of total EU GDP through consumption and that its real spending power has been growing faster than the rest of the population's for the last decade. In Italy alone, the over-65 segment commanded approximately \texteuro{}243.8~billion in aggregate annual household income flow and approximately \texteuro{}2.33~trillion in household net wealth in 2024, computed as Istat over-65 individual counts multiplied by SHARE Wave~9 median household figures.

This is the silver economy. It is the largest customer base of contemporary Italian capitalism, and unlike younger demographic blocks it is not declining: the segment that funds private pensions, owns the housing stock, sits on the financial assets, and consumes a disproportionate share of healthcare, financial advisory and personal-services markets is structurally expanding through 2050. For a country whose macroeconomic growth is constrained by labour-force decline, the productivity of the silver-economy segment --- the efficiency with which markets serve it --- is increasingly indistinguishable from the productivity of the economy as a whole.

\section{The system treats the segment as a monolith. It is not.}
Despite the segment's economic significance, Italian commercial and public services continue to address the over-65 population as if it were a single homogeneous block defined by age. Insurance companies design "senior" products tailored to a stylised retiree; banks deploy generic "wealth preservation" packages; the public administration plans elderly services from aggregate municipal head-counts; pharmaceutical and OTC operators run age-band marketing campaigns that ignore the substantial heterogeneity of digital reach and engagement within the cohort. The premise of this thesis is that this monolithic treatment is both behaviourally wrong and economically inefficient.

The heterogeneity is empirically large. Among Italians aged 65 and over, internet penetration in 2021--2022 ranged from 9\% in the most isolated profile to 75\% in the most digitally engaged; financial distress affected roughly half of one segment and almost no respondents in another; quality of life as measured by the CASP-12 scale spanned more than ten points across profiles within the same country, age group and clinical condition. A "senior" insurance product positioned for the average over-65 customer is a product positioned for nobody in particular: it under-serves the digital high-income segment that wants self-managed online onboarding, fails the asset-rich socially-isolated segment that needs human-mediated navigation, and prices uncompetitively for the fragile lower-income segment whose actual willingness-to-pay is below the average premium.

The cost of this misalignment is observable in two ways that this thesis quantifies in subsequent chapters. The first is healthcare under-utilisation. Italy's over-65 population pays for routine general-practitioner contact at a rate comparable to Sweden's but uses preventive services --- dental care, elective specialist consultation, screening --- at a fraction of the Swedish rate (Italy 30\% twelve-month dentist contact against Sweden 84\%). Within Italy, this preventive-care deficit concentrates in a profile whose objective conditions are good and whose ability to pay is intact, but whose social and digital integration is structurally low. Chapter~\ref{ch:ch7} documents that this Moderate Isolated profile uses specialist services at half the rate of an otherwise comparable Traditional Social profile, yet does not report any forgone-care-for-cost gap; the deficit is one of \emph{navigational capital}, not of affordability.

The second is the digital transition. Italy's National Recovery and Resilience Plan (PNRR) Mission 6 invests \texteuro{}15.63~billion in digital healthcare infrastructure under the assumption that digital channels improve access for all citizens. For the segment of the over-65 population that has the income to pay for premium services, the digital infrastructure works as intended; for the segment that lacks digital reach, the same infrastructure functions as a mechanism of selective exclusion. Without a segmentation that lets policymakers and operators see the heterogeneity ex~ante, the digital transition risks reproducing inequality at scale.

\section{Sweden as the operational benchmark}
The contrast with Sweden is a useful one because Sweden is, in a precise sense, the most operationally mature elderly-services market in Europe. The Swedish welfare regime has historically allocated approximately three times as large a share of GDP to long-term care as the Italian regime, but this institutional asymmetry has secondary effects on the private market that are at least as relevant for the silver-economy framing of this thesis: Swedish elderly consumers are more digitally engaged (88.8\% internet penetration in the SHARE W9 sample versus 39.0\% in Italy), more accustomed to preventive services, more financially mobile, and consequently more receptive to a richer set of private products in financial services, real estate, retail and digital health than their Italian peers. Where Swedish private operators have already built scaled offerings --- private dental insurance, hybrid digital-physical wealth management, senior co-housing, on-demand care subscriptions --- these markets serve as observable references for what a more operationally mature Italian market could look like.

The thesis does not argue that Italy should converge on the Swedish welfare model. It argues that, profile by profile, the Italy--Sweden gap quantifies an addressable headroom for Italian operators and that the gap concentrates in segments whose Italian counterparts are demonstrably solvent. Sweden in this thesis is a reference point for adoption curves and product viability, not a normative ideal.

\section{A segmentation that produces operational decisions}
The segmentation framework deployed in this thesis operationalises the silver-economy view of the over-65 population on five conceptual dimensions: health, economic resources, digital and social engagement, cognitive functioning, and subjective wellbeing. The five-dimensional framework is deliberate. The standard segmentations available to commercial operators are typically two- or three-dimensional --- variations of clinical risk, income band and household composition --- and produce profiles that perform poorly on the discretionary-consumption questions that matter for the silver-economy decisions. Adding cognitive functioning makes the segment behaviourally legible (the literature on cognitive reserve, e.g.\ Stern, 2002, ties cognition to engagement with complex services). Adding subjective wellbeing as a clustering input rather than as a downstream validation outcome separates two profiles --- Fragile Resigned and Fragile Depressed --- that are clinically indistinguishable but commercially distinct: the Resigned profile is an outreach problem, the Depressed profile is a caregiver-mediated channel problem.

The methodological commitment to subjective inputs is non-trivial and is one of the three original contributions of the thesis (see Chapter~\ref{ch:ch2}). The empirical payoff is documented in Chapter~\ref{ch:ch6}: removing the subjective block from the clustering input collapses a substantively important distinction, reduces the external ANOVA $F$-statistic on life satisfaction by 41\% in Italy and 22\% in Sweden, and produces a partition that disagrees with the five-dimensional reference on roughly half of the assignments.

\section{Research questions}
Three interrelated questions organise the empirical work of this thesis.

\textbf{RQ1.} How does the Italian over-65 population segment when measured on five dimensions (health, economic resources, digital and social engagement, cognitive functioning, subjective wellbeing), which behaviourally and economically distinct profiles emerge, and how large are they in addressable individuals and \texteuro{}?

\textbf{RQ2.} How does the Italian profile structure compare with Sweden --- the most operationally mature European silver-economy market --- and where do the matched-pair gaps quantify operational headroom for Italian private operators and policymakers?

\textbf{RQ3.} How can the segmentation be translated into a deployable decision-support instrument that allows commercial operators (insurance, wealth management, senior services, pharma marketing) and public administration to make profile-aware investment, product, and policy decisions?

RQ1 is addressed in Chapter~\ref{ch:ch4}, which presents the five-profile Italian segmentation, with profile sizes ranging from 184 (Traditional Social, 7.7\% of the sample) to 639 (Moderate Isolated, 26.9\%) and aggregate annual income flow per profile ranging from approximately \texteuro{}27~billion to \texteuro{}96~billion. RQ2 is addressed in Chapter~\ref{ch:ch5}, which produces a six-profile Swedish segmentation under the identical analytical pipeline, identifies four matched pairs through a semantic mapping documented in Chapter~\ref{ch:ch3}, and quantifies the Sweden-minus-Italy gap on the dimensions that map onto market headroom (CASP, internet, dental and specialist contact, financial ease). RQ3 is addressed in Chapter~\ref{ch:ch8}, which translates the framework into the Italian Silver Atlas --- a publicly deployed decision-support tool with four screens (Overview, Atlas, Benchmark, Opportunity Explorer) and four vertical playbooks (private health insurance, wealth management, senior living and home services, pharma and OTC marketing).

\section{Three original contributions}
The thesis makes three original contributions to the literature and practice of elderly segmentation.

The first is methodological. Subjective wellbeing measures are incorporated as inputs to the segmentation rather than as validation outcomes, departing from the established SHARE segmentation literature (Lestari et~al., 2022; Rodrigues et~al., 2023; Han et~al., 2025). The empirical justification is the 4D-versus-5D sensitivity analysis of Chapter~\ref{ch:ch6}, which establishes that removing the subjective block both reduces external discriminative power on the life-satisfaction outcome and collapses the Fragile Resigned / Fragile Depressed distinction.

The second is comparative. The thesis conducts a profile-level comparison between a Southern-European (Italy) and a Northern-European (Sweden) over-65 population using identical variables and methods, with the explicit goal of identifying matched pairs and welfare-regime-specific profiles. To the best of this author's knowledge, no prior study has performed a five-dimensional comparative segmentation across these two welfare regimes using SHARE Wave~9 data, and none has framed the resulting gap as addressable market headroom for Italian operators.

The third is operational. The thesis instantiates the segmentation as a publicly deployed product --- the Italian Silver Atlas --- that is usable in autonomy by a strategy analyst, a product manager, an investor, or a public-administration planner. The product is not a demonstrative appendix to the academic work; it is the principal deliverable of the thesis under the silver-economy framing, and it is designed to support recurring real-world decisions: investment go/no-go on a vertical, customer segmentation in a CRM, board-memo sizing for a new product, white-space identification against the Swedish benchmark.

\section{Structure of the thesis}
The remainder of the thesis is organised as follows. Chapter~\ref{ch:ch2} reviews the literature on elderly segmentation, the digital divide, the objective-versus-subjective distinction in wellbeing measurement, and the silver-economy framing. Chapter~\ref{ch:ch3} describes the data, the variables and the methodological pipeline (PCA, factor analysis, K-means and Latent Class Analysis). Chapter~\ref{ch:ch4} presents the Italian segmentation. Chapter~\ref{ch:ch5} presents the cross-country comparison with Sweden. Chapter~\ref{ch:ch6} documents the robustness and sensitivity analyses, including the methodologically central 4D-versus-5D sensitivity. Chapter~\ref{ch:ch7} examines healthcare-utilisation patterns across profiles, focusing on the Moderate Isolated--Traditional Social contrast and reframing it as a navigational-capital deficit. Chapter~\ref{ch:ch8} describes the Italian Silver Atlas product, articulates the go-to-market and the four vertical playbooks, and concludes.
